
\begin{frame}{Appendix --- Theoretical Framework: Signal vs Weapon}
\hypertarget{app:theory}{}
\footnotesize
The \textbf{Leveling} and \textbf{Barrier} faces of the main text rest on two hypotheses about \emph{why} litigation moves an election. A third channel --- \emph{direct judicial action}, where a ruling removes a candidate and mechanically reallocates the vote --- lies \textbf{outside the estimand}: the adversarial filter drops the candidacy-registration and campaign-finance machinery (RRC/DRAP, \emph{presta\c{c}\~ao de contas}), so it is not what these estimates measure.
\smallskip
\begin{columns}[T]
\begin{column}{0.49\textwidth}
  \begin{block}{H1 --- Accountability / signal}
  \scriptsize Courts sanction abuse; a conviction disclosed before the vote reveals candidate type $\Rightarrow$ cleaner competitors advance. \hfill\ns{\citep{assumpcao2019}}
  \end{block}
\end{column}
\begin{column}{0.49\textwidth}
  \begin{block}{H2 --- Weaponization / noise}
  \scriptsize Litigation as campaign weapon: aimed up the ladder, judicially weak, featured in media campaigns $\Rightarrow$ aimed at voters, not judges. \hfill\ns{\citep{chin2025,nakaguma2025}}
  \end{block}
\end{column}
\end{columns}
\smallskip

The testable split is Leveling against Barrier rather than H1 against H2. Voters cannot tell a genuine charge from a strategic one, so an accountability signal (H1) and a \emph{sticky} weapon (an H2 that changes minds) look alike: both move votes toward someone, which is the \textbf{Leveling} signature. What the data can separate is whether litigation moves votes at all (Leveling) or degrades the contest without re-sorting (Barrier). The three mechanisms collapse onto one observable axis:
\smallskip
\begin{center}\scriptsize
\renewcommand{\arraystretch}{1.3}
\begin{tabular}{ll}
\toprule
\textbf{Mechanism} & \textbf{Observable face}\\
\midrule
H1 accountability signal & \textbf{Leveling} --- cleaner types advance; votes re-sort\\
H2 weapon that \emph{sticks} (persuades) & \textbf{Leveling} --- votes move; indistinguishable from H1\\
H2 weapon that only \emph{fouls} & \textbf{Barrier} --- blank/null $\uparrow$, valid $\downarrow$; no re-sorting\\
\bottomrule
\end{tabular}
\end{center}
\smallskip
\centerline{The design identifies Leveling against Barrier; it does not separate H1 from a sticky H2. \hfill\hyperlink{src:motivation}{\beamerreturnbutton{Back}}}
\end{frame}
